\section{Kinematics}\label{sec:kinematics}

\subsection{Displacement, Velocity, and Acceleration}\label{sec:displacement_velocity_acceleration}

I will avoid using SUVAT equations immediately, since they hide a lot of the fundamentals.

We can relate displacement, velocity, and acceleration as functions of time $t$ via integration and differentiation. This is by recalling that 
velocity is the rate of change of displacement, and acceleration is the rate of change of velocity, so we can write:

\begin{align*}
    v(t) &= \frac{ds}{dt} \\
    a(t) &= \frac{dv}{dt} = \frac{d^2s}{dt^2}
\end{align*}

And you can go the `other way' via integration:
\begin{align*}
    v(t) &= \int a(t) dt + c_1 \\
    s(t) &= \int v(t) dt + c_2
\end{align*}

Where we can find the constants of integration $c_1$ and $c_2$ if we know some initial conditions, e.g. the starting velocity and the starting position.

Let's say we have a particle moving with \textbf{constant acceleration} $a$, with an initial velocity $u$ and initial displacement $s_0$. 

We can start with the acceleration and integrate\footnote{It is essential to understand that this integral relies on $a$ being constant in $t$.} to find:
\[ v(t) = \int a dt + c  = at + c_1\]
And use that $v(0) = u$ to find:
\begin{align*}
    v(0) &= u \\
    &=  a \cdot 0 + c_1  \\
    \therefore c_1 &= u\\
    \therefore v(t) &= at + u
\end{align*}


Now we can integrate again to find the displacement:
\[ s(t) = \int v(t) dt + c_2 = \int (at + u) dt + c_2 = \frac{1}{2}at^2 + ut + c_2 \]

And use that $s(0) = s_0$ to find:
\begin{align*}
    s(0) &= s_0 \\
    &= \frac{1}{2}a \cdot 0^2 + u \cdot 0 + c_2 \\
    \therefore c_2 &= s_0 \\
    \therefore s(t) &= \frac{1}{2}at^2 + ut + s_0
    \therefore s(t) - s_0 = ut + \frac{1}{2}at^2
\end{align*}

We can write this in terms of $v$ instead of $u$ by substituting the first equation $v = at + u$ to get:
\begin{align*}
    s(t) - s_0 &= ut + \frac{1}{2}at^2 \\
    &= (v - at)t + \frac{1}{2}at^2 \\
    &= vt - at^2 + \frac{1}{2}at^2 \\
    &= vt - \frac{1}{2}at^2
\end{align*}

Now we can notice that the last two equations have a similar term $\pm \frac{1}{2}at^2$, so we can add them together to get:
\begin{align*}
    2(s(t) - s_0) &= ut + vt  + \cancel{\frac{1}{2}at^2 - \frac{1}{2}at^2}\\
    2(s(t) - s_0) &= t(u + v) \\
    \therefore s(t) - s_0 &= \frac{t(u + v)}{2}
\end{align*}

All of the last equations contain $t$ so it would be useful to eliminate $t$\footnote{The neatest way to do this is via chain rule and writing $a = \frac{dv}{dt} = \frac{dv}{ds} \cdot \frac{ds}{dt} = v \frac{dv}{ds}$, but we can also do it via substitution, which I suspect students will find simpler to understand.}. 
If we assume $a \neq 0$, we can rearrange the first equation to get $t = \frac{v - u}{a}$, and substitute this into the last equation to get:
\begin{align*}
    s(t) - s_0 &= \frac{t(u + v)}{2} \\
    &= \frac{\frac{v - u}{a}(u + v)}{2} \\
    &= \frac{(v - u)(u + v)}{2a} \\
    &= \frac{v^2 - u^2}{2a}
    \therefore v^2&= u^2 + 2a(s - s_0)
\end{align*}

These combine to give us our SUVAT equations from the formula book, explaining motion in a straight line with constant acceleration $a$, with initial velocity $u$, initial displacement $s_0$, and displacement $s(t)$ and velocity $v(t)$ at time $t$:

\begin{align*}
    v(t) &= u + at \\
    s(t) - s_0 &= ut + \frac{1}{2}at^2 \\
    s(t) - s_0 &= v(t)t - \frac{1}{2}at^2 \\
    s(t) - s_0 &= \frac{t(u + v(t))}{2} \\
    v(t)^2&= u^2 + 2a(s(t) - s_0)
\end{align*}

In particular, if $s_0 = 0$\footnote{This highlights that the SUVAT equations describe displacement from the starting position, rather than absolute position relative to a fixed origin.}, i.e. we are measuring relative to the starting point, and we set $v$ as our final velocity and $s$ as our displacement (ignoring the time variable), then we get the SUVAT equations from the formula booklet:
\begin{align*}
    v &= u + at \\
    s &= ut + \frac{1}{2}at^2 \\
    s &= vt - \frac{1}{2}at^2 \\
    s &= \frac{t(u + v)}{2} \\
    v^2&= u^2 + 2as
\end{align*}

I highly recommend understanding where the equations come from, especially relating to differentiating and integrating acceleration, velocity, and displacement, rather than just blindly using the equations.

Despite this, the equations can still speed up calculations, since they always allow you to find the missing variable without needing to do any differentiation or integration. 
This is because each equation has $4$ of the variables $s$, $u$, $v$, $a$, and $t$, so if you know any $3$ of these variables, you can find the missing variable by using the appropriate SUVAT equation.

The only place you might get tripped up is forgetting that $s$ is relative to the starting point, in the case where you are not starting from $s_0 = 0$.

\begin{example}
A particle starts from rest and moves with constant acceleration $a$ for $t$ seconds, covering a distance of $5$ metres. 

Find an equation for the final velocity $v$ of the particle in terms of $s$ and $t$, assuming $t \neq 0$.

You are told the particle reaches a final velocity of $2$ m/s. Find what time $t$ is at the time of reaching the final velocity, and then deduce what the acceleration $a$ is.

\end{example}
\begin{solution}
Since the particle starts from rest, we have $u = 0$. We are asked to find an equation for $v$ in terms of $s$ and $t$ and we know what $u$ is, so let's look for a SUVAT equation which has $v$, $s$, $t$, and $u$. The only equation which has all of these variables is the fourth equation, so we can rearrange this to get:
\begin{align*}
    s &= \frac{t(u + v)}{2} \\
    2s &= t(u + v) \\
    2s &= t(0 + v) \\
    2s &= tv \\
    \therefore v &= \frac{2s}{t} \text{ m/s}
\end{align*}
We are told that $v = 2$ m/s after the object has covered $5$ metres, so we can substitute this into the equation we just found to get:
\begin{align*}
    2 &= \frac{2 \cdot 5}{t} \\
    t &= 5 \text{ seconds}
\end{align*}
Now we can use the first SUVAT equation to find $a$:
\begin{align*}
    v &= u + at \\
    2 &= 0 + a \cdot 5 \\
    a &= \frac{2}{5} \text{ m/s}^2
\end{align*}
\qed 
\end{solution}

\begin{example}
A particle starts with velocity $6$ m/s and moves with constant acceleration $a$, covering a distance of $32$ metres before reaching a velocity $v$.

Find an equation for the acceleration $a$ in terms of $v$.

You are told the particle reaches a final velocity of $10$ m/s. Find the acceleration $a$, and then deduce the time $t$ taken.

\end{example}
\begin{solution}
We know that $u = 6$ and $s = 32$. We are asked to find an equation for $a$ in terms of $v$, so we look for a SUVAT equation involving $a$, $v$, $u$, and $s$. The only equation with these variables is:
\begin{align*}
    v^2 &= u^2 + 2as
\end{align*}
Rearranging this to make $a$ the subject gives:
\begin{align*}
    v^2 &= u^2 + 2as \\
    v^2 - u^2 &= 2as \\
    a &= \frac{v^2 - u^2}{2s}
\end{align*}
Substituting $u = 6$ and $s = 32$, we get:
\begin{align*}
    a &= \frac{v^2 - 6^2}{2 \cdot 32} \\
    &= \frac{v^2 - 36}{64}
\end{align*}
So the required equation is:
\[
a = \frac{v^2 - 36}{64} \text{ m/s}^2
\]

We are told that $v = 10$ m/s, so:
\begin{align*}
    a &= \frac{10^2 - 36}{64} \\
    &= \frac{100 - 36}{64} \\
    &= \frac{64}{64} \\
    &= 1 \text{ m/s}^2
\end{align*}

Now we use the first SUVAT equation to find the time:
\begin{align*}
    v &= u + at \\
    10 &= 6 + 1 \cdot t \\
    t &= 4 \text{ seconds}
\end{align*}

Hence the five variables are:
\[
u = 6,\quad v = 10,\quad s = 32,\quad a = 1,\quad t = 4
\]
\qed
\end{solution}

\begin{example}
A particle moves with constant acceleration $a$, changing its velocity from $2$ m/s to $14$ m/s in $6$ seconds, while travelling a distance $s$.

Find an equation for the displacement $s$ in terms of $u$.

You are told that the initial velocity is $2$ m/s. Find the displacement $s$, and then deduce the acceleration $a$.

\end{example}
\begin{solution}
We know that $v = 14$ and $t = 6$. We are asked to find an equation for $s$ in terms of $u$, so we look for a SUVAT equation involving $s$, $u$, $v$, and $t$. The only equation with these variables is:
\begin{align*}
    s &= \frac{t(u + v)}{2}
\end{align*}
Substituting $v = 14$ and $t = 6$, we get:
\begin{align*}
    s &= \frac{6(u + 14)}{2} \\
    &= 3(u + 14)
\end{align*}
So the required equation is:
\[
s = 3(u + 14) \text{ metres}
\]

We are told that $u = 2$ m/s, so:
\begin{align*}
    s &= 3(2 + 14) \\
    &= 3 \cdot 16 \\
    &= 48 \text{ metres}
\end{align*}

Now we use the first SUVAT equation to find the acceleration:
\begin{align*}
    v &= u + at \\
    14 &= 2 + a \cdot 6 \\
    12 &= 6a \\
    a &= 2 \text{ m/s}^2
\end{align*}

Hence the five variables are:
\[
u = 2,\quad v = 14,\quad t = 6,\quad s = 48,\quad a = 2
\]
\qed
\end{solution}

\begin{example}
A particle starts at position $s_0 = 4$ metres with initial velocity $u = 3$ m/s and moves with constant acceleration $a = 2$ m/s$^2$ for $t$ seconds, reaching position $s$.

Find the position $s$ after $5$ seconds.

\end{example}
\begin{solution}
Since the particle does not start at the origin, we must be careful to use the SUVAT equation for displacement relative to the starting position:
\begin{align*}
    s - s_0 &= ut + \frac{1}{2}at^2
\end{align*}
We are given:
\[
s_0 = 4,\quad u = 3,\quad a = 2,\quad t = 5
\]
So substituting into the equation gives:
\begin{align*}
    s - 4 &= 3 \cdot 5 + \frac{1}{2}\cdot 2 \cdot 5^2 \\
    s - 4 &= 15 + 25 \\
    s - 4 &= 40 \\
    s &= 44
\end{align*}
Therefore the position of the particle after $5$ seconds is:
\[
s = 44 \text{ metres}
\]
\qed
\end{solution}

\begin{example}
    A particle at rest which has a mass of $2$kg has a force applied to it of $4$N. 

    What is the acceleration of the particle?

    Deduce an equation for the velocity $v$ of the particle in terms of time $t$.
\end{example}
\begin{solution}
We know that the force $F$ and mass $m$ are related to acceleration $a$ by Newton's second law:
\begin{align*}
    F &= ma
\end{align*}
Substituting $F = 4$ N and $m = 2$ kg, we get:
\begin{align*}
    4 &= 2a \\
    a &= 2 \text{ m/s}^2
\end{align*}

Since the particle starts from rest, the initial velocity $u = 0$. Using the first SUVAT equation:
\begin{align*}
    v &= u + at \\
    v &= 0 + 2t \\
    v &= 2t
\end{align*}
So the velocity of the particle in terms of time $t$ is:
\[
v = 2t \text{ m/s}
\]
\qed
\end{solution}

\begin{example}
    A particle has a mass of $2$kg and has been moving at a constant speed of $3$ m/s. At time $t=0$ seconds, a force of $4t$N is applied opposing the motion.

    What is the acceleration of the particle in terms of time $t$?

    Deduce an equation for the velocity $v$ of the particle in terms of time $t$.

    At what time $t$ does the particle come to rest?
\end{example}
\begin{solution}
Let's set our positive direction to be in the direction of the initial motion.Using Newton's second law, we have:
\begin{align*}
    F &= ma \\
    -4t &= 2a \\
    a &= -2t \text{ m/s}^2
\end{align*}
Since the particle is moving at a constant speed of $3$ m/s before the force is applied, we have $u = 3$. 

Now the question is trying to trick us; we do not have constant acceleration here. Instead we must remember where the equations came from:

\[ v(t) = \int a(t) \, dt + c\]

Substituting $a(t) = -2t$ gives:
\begin{align*}
    v(t) &= \int -2t \, dt + c \\
    &= -t^2 + c
\end{align*}
Since the initial velocity is $u = 3$ m/s, we have $v(0) = 3$, so $c = 3$. Therefore:
\begin{align*}
    v(t) &= -t^2 + 3
\end{align*}

The particle comes to rest when $v(t) = 0$, so we can solve the equation:
\begin{align*}
    0 &= -t^2 + 3 \\
    t^2 &= 3 \\
    t &= \sqrt{3} \text{ seconds}
\end{align*}

\qed
\end{solution}


\begin{example}
    A star is being consumed by a black hole at a large distance. The star is losing mass, and is described by $M/t$ kg where $M$ is the initial mass of the star. The black hole is exerting a force of $100t^2$ N on the star, and the star is accelerating towards the black hole.

    Modelling the star as a particle, find an equation for the acceleration $a$ of the star in terms of time $t$.

    Assuming the star starts from rest, find an equation for the velocity $v$ of the star in terms of time $t$.
\end{example}

\begin{solution}
Using Newton's second law, we have:
\begin{align*}
    F &= ma \\
    100t^2 &= \frac{M}{t}a \\
    a &= \frac{100t^2 \cdot t}{M} \\
    a &= \frac{100t^3}{M} \text{ m/s}^2
\end{align*}

Since the star starts from rest, the initial velocity $u = 0$. Using the definition of acceleration:
\begin{align*}
    v(t) &= \int a(t) \, dt + c \\
    &= \int\frac{100t^3}{M} \, dt + c \\
     &= \frac{100t^4}{4M} + c \\
     &= \frac{25t^4}{M} + c
\end{align*}
Since $v(0) = 0$, we have 
\begin{align*}
    0 &= \frac{25 \cdot 0^4}{M} + c \\
    \therefore c &= 0
\end{align*}
Therefore:
\begin{align*}
    v(t) &= \frac{25t^4}{M}
\end{align*}
\qed

\end{solution}